\begin{thebibliography}{99}
\bibitem{bib-bib1} 
\par\noindent    Bahareh Afshari and Graham E. Leigh.  Lyndon interpolation for modal $\mu$-calculus.  In Aybüke Özgün and Yulia Zinova, editors, Language, Logic, and Computation - 13th International Tbilisi Symposium, TbiLLC 2019, Batumi, Georgia, September 16-20, 2019, Revised Selected Papers, volume 13206 of Lecture Notes in Computer Science, pages 197–213. Springer, 2019.  doi:10.1007/978-3-030-98479-3\_10.  

\bibitem{bib-bib2} 
\par\noindent    Michele Basaldella and Kazushige Terui.  Infinitary completeness in ludics.  In Proceedings of the 25th Annual IEEE Symposium on Logic in Computer Science, LICS 2010, 11-14 July 2010, Edinburgh, United Kingdom, pages 294–303. IEEE Computer Society, 2010.  doi:10.1109/LICS.2010.47.  

\bibitem{bib-bib3} 
\par\noindent    Michele Basaldella and Kazushige Terui.  On the meaning of logical completeness.  Log. Methods Comput. Sci., 6(4), 2010.  doi:10.2168/LMCS-6(4:11)2010.  

\bibitem{bib-bib4} 
\par\noindent    Evert W Beth.  On padoa’s method in the theory of definition.  Indagationes Mathematicae, 15:330–339, 1953.  

\bibitem{bib-bib5} 
\par\noindent    James Brotherston and Rajeev Goré.  Craig interpolation in displayable logics.  In TABLEAUX, volume 6793 of Lecture Notes in Computer Science, pages 88–103. Springer, 2011.  doi:10.1007/978-3-642-22119-4\_9.  

\bibitem{bib-bib6} 
\par\noindent    Alessandra Carbone.  Interpolants, cut elimination and flow graphs for the propositional calculus.  Ann. Pure Appl. Log., 83(3):249–299, 1997.  doi:10.1016/S0168-0072(96)00019-X.  

\bibitem{bib-bib7} 
\par\noindent    William Craig.  Linear reasoning. A new form of the Herbrand-Gentzen theorem.  J. Symb. Log., 22(3):250–268, 1957.  doi:10.2307/2963593.  

\bibitem{bib-bib8} 
\par\noindent    William Craig.  Three uses of the herbrand-gentzen theorem in relating model theory and proof theory.  J. Symb. Log., 22(3):269–285, 1957.  doi:10.2307/2963594.  

\bibitem{bib-bib9} 
\par\noindent    Djordje Cubric.  Results in categorical proof theory.  PhD thesis, McGill University, 1993.  

\bibitem{bib-bib10} 
\par\noindent    Djordje Cubric.  Interpolation property for bicartesian closed categories.  Arch. Math. Log., 33(4):291–319, 1994.  doi:10.1007/BF01270628.  

\bibitem{bib-bib11} 
\par\noindent    Pierre-Louis Curien and Hugo Herbelin.  The duality of computation.  In Martin Odersky and Philip Wadler, editors, Proceedings of the Fifth ACM SIGPLAN International Conference on Functional Programming (ICFP ’00), Montreal, Canada, September 18-21, 2000, pages 233–243. ACM, 2000.  doi:10.1145/351240.351262.  

\bibitem{bib-bib12} 
\par\noindent    Pierre-Louis Curien and Guillaume Munch-Maccagnoni.  The duality of computation under focus.  In Cristian S. Calude and Vladimiro Sassone, editors, Theoretical Computer Science - 6th IFIP TC 1/WG 2.2 International Conference, TCS 2010, Held as Part of WCC 2010, Brisbane, Australia, September 20-23, 2010. Proceedings, volume 323 of IFIP Advances in Information and Communication Technology, pages 165–181. Springer, 2010.  doi:10.1007/978-3-642-15240-5\_13.  

\bibitem{bib-bib13} 
\par\noindent    Vincent Danos, Jean-Baptiste Joinet, and Harold Schellinx.  A new deconstructive logic: Linear logic.  J. Symb. Log., 62(3):755–807, 1997.  doi:10.2307/2275572.  

\bibitem{bib-bib14} 
\par\noindent    Jeremy E. Dawson, James Brotherston, and Rajeev Goré.  Machine-checked interpolation theorems for substructural logics using display calculi.  In IJCAR, volume 9706 of Lecture Notes in Computer Science, pages 452–468. Springer, 2016.  doi:10.1007/978-3-319-40229-1\_31.  

\bibitem{bib-bib15} 
\par\noindent    Adrien Deloro.  Cours de logique mathématique.  Draft of a book on mathematical logic, 2025.  URL: https://webusers.imj-prg.fr/˜adrien.deloro/index.php?topic=teaching\&page=logique.  

\bibitem{bib-bib16} 
\par\noindent    Kosta Došen.  Logical Consequence: A Turn in Style, pages 289–311.  Springer Netherlands, Dordrecht, 1997.  doi:10.1007/978-94-017-0487-8\_15.  

\bibitem{bib-bib17} 
\par\noindent    Solomon Feferman.  Harmonious logic: Craig’s interpolation theorem and its descendants.  Synthèse, 164(3):341–357, 2008.  doi:10.1007/S11229-008-9354-2.  

\bibitem{bib-bib18} 
\par\noindent    Melvin Fitting.  Linear reasoning in modal logic.  J. Symb. Log., 49(4):1363–1378, 1984.  doi:10.2307/2274285.  

\bibitem{bib-bib19} 
\par\noindent    Jean-Yves Girard.  Linear logic.  Theoretical Computer Science, 50(1):1–101, 1987.  doi:10.1016/0304-3975(87)90045-4.  

\bibitem{bib-bib20} 
\par\noindent    Jean-Yves Girard.  Proof theory and logical complexity.  Bibliopolis, 1987.  

\bibitem{bib-bib21} 
\par\noindent    Thomas A. Henzinger, Ranjit Jhala, Rupak Majumdar, and Kenneth L. McMillan.  Abstractions from proofs.  SIGPLAN Not., 39(1):232–244, January 2004.  doi:10.1145/982962.964021.  

\bibitem{bib-bib22} 
\par\noindent    Stefan Hetzl and Raheleh Jalali.  On the completeness of interpolation algorithms.  In Pawel Sobocinski, Ugo Dal Lago, and Javier Esparza, editors, Proceedings of the 39th Annual ACM/IEEE Symposium on Logic in Computer Science, LICS 2024, Tallinn, Estonia, July 8-11, 2024, pages 45:1–45:13. ACM, 2024.  doi:10.1145/3661814.3662112.  

\bibitem{bib-bib23} 
\par\noindent    Makoto Kanazawa.  Computing interpolants in implicational logics.  Annals of Pure and Applied Logic, 142(1):125–201, 2006.  doi:10.1016/j.apal.2005.12.014.  

\bibitem{bib-bib24} 
\par\noindent    Roman Kuznets.  Multicomponent proof-theoretic method for proving interpolation properties.  Ann. Pure Appl. Log., 169(12):1369–1418, 2018.  doi:10.1016/J.APAL.2018.08.007.  

\bibitem{bib-bib25} 
\par\noindent    Shoji Maehara.  On the interpolation theorem of Craig.  Sûgaku, 12(4), 1960.  

\bibitem{bib-bib26} 
\par\noindent    Ralph Matthes.  Interpolation for natural deduction with generalized eliminations.  In Reinhard Kahle, Peter Schroeder-Heister, and Robert Stärk, editors, Proof Theory in Computer Science, pages 153–169, Berlin, Heidelberg, 2001. Springer Berlin Heidelberg.  doi:10.1007/3-540-45504-3\_10.  

\bibitem{bib-bib27} 
\par\noindent    Kenneth L. McMillan.  Interpolation and model checking.  In Edmund M. Clarke, Thomas A. Henzinger, Helmut Veith, and Roderick Bloem, editors, Handbook of Model Checking., pages 421–446. Springer, 2018.  doi:10.1007/978-3-319-10575-8\_14.  

\bibitem{bib-bib28} 
\par\noindent    Grigori Mints.  Finite investigations of transfinite derivations.  Journal of Soviet Mathematics, 10:548–596, 1978.  

\bibitem{bib-bib29} 
\par\noindent    Guillaume Munch-Maccagnoni.  Focalisation and Classical Realisability.  In Erich Grädel and Reinhard Kahle, editors, Computer Science Logic ’09, volume 5771 of Lecture Notes in Computer Science, pages 409–423. Springer, Heidelberg, 2009.  doi:10.1007/978-3-642-04027-6\_30.  

\bibitem{bib-bib30} 
\par\noindent    Guillaume Munch-Maccagnoni.  Syntax and Models of a non-Associative Composition of Programs and Proofs.  PhD thesis, Univ. Paris Diderot, 2013.  

\bibitem{bib-bib31} 
\par\noindent    Dag Prawitz.  Natural Deduction: A Proof-Theoretical Study.  Dover Publications, Mineola, N.Y., 1965.  

\bibitem{bib-bib32} 
\par\noindent    João Rasga, Walter Alexandre Carnielli, and Cristina Sernadas.  Interpolation via translations.  Math. Log. Q., 55(5):515–534, 2009.  doi:10.1002/MALQ.200810013.  

\bibitem{bib-bib33} 
\par\noindent    A. Robinson.  A result of consistency and its application to the theory of definition.  Indagationes Mathematicae, 18, 1956.  

\bibitem{bib-bib34} 
\par\noindent    Dirk Roorda.  Interpolation in fragments of classical linear logic.  The Journal of Symbolic Logic, 59(2):419–444, 1994.  doi:10.2307/2275398.  

\bibitem{bib-bib35} 
\par\noindent    Kurt Schütte.  Proof theory, volume 85.  Springer, 1977.  

\bibitem{bib-bib36} 
\par\noindent    D. S. Shamkanov.  Circular proofs for the Gödel-Löb provability logic.  Mathematical Notes, 96(3), 2014.  doi:10.1134/S0001434614090326.  

\bibitem{bib-bib37} 
\par\noindent    Lutz Strassburger.  Linear Logic and noncommutativity in the Calculus of Structures.  PhD thesis, TU Dresden, 2003.  

\bibitem{bib-bib38} 
\par\noindent    Amirhossein Akbar Tabatabai and Raheleh Jalali.  Universal proof theory: Semi-analytic rules and craig interpolation.  Ann. Pure Appl. Log., 176(1):103509, 2025.  doi:10.1016/J.APAL.2024.103509.  

\bibitem{bib-bib39} 
\par\noindent    Gaisi Takeuti.  Proof theory, volume 81.  Courier Corporation, 2013.  

\bibitem{bib-bib40} 
\par\noindent    A. S. Troelstra and H. Schwichtenberg.  Basic Proof Theory.  Cambridge Tracts in Theoretical Computer Science. Cambridge University Press, 2 edition, 2000.  

\end{thebibliography}
